En el anexo se recomienda dejar la evidencia completa del desarrollo computacional, sin saturar el cuerpo principal del informe. Los archivos generados por el script que respaldan el an\'alisis son, entre otros, \texttt{boxplots	extunderscore outliers.pdf}, \texttt{scree	extunderscore plot.pdf}, \texttt{relacion	extunderscore rm	extunderscore medv.pdf}, \texttt{boxplot	extunderscore medv	extunderscore chas.pdf}, \texttt{relacion	extunderscore lstat	extunderscore medv.pdf}, \texttt{regresion	extunderscore medv	extunderscore rm.txt}, \texttt{regresion	extunderscore medv	extunderscore lstat.txt}, \texttt{regresion	extunderscore multiple	extunderscore crim.txt} y \texttt{respuestas	extunderscore preguntas.txt}.

\subsection{Bit\'acora de trabajo}
\begin{itemize}
    \item D\'ia 1: revisi\'on del enunciado, carga del dataset Boston y exploraci\'on inicial.
    \item D\'ia 2: limpieza, estandarizaci\'on, PCA y generaci\'on de figuras/tablas.
    \item D\'ia 3: interpretaci\'on de resultados, redacci\'on del informe y validaci\'on final.
\end{itemize}